\documentclass[11pt]{article}

% --- Conservative preamble: only standard, arXiv-safe packages ---
\usepackage[utf8]{inputenc}
\usepackage[T1]{fontenc}
\usepackage{amsmath, amssymb, amsthm}
\usepackage{geometry}
\geometry{margin=1in}
\usepackage{hyperref}
\hypersetup{colorlinks=true, linkcolor=blue, citecolor=blue, urlcolor=blue}
\usepackage{booktabs}
\usepackage{array}

\newtheorem{theorem}{Theorem}
\newtheorem{proposition}{Proposition}
\newtheorem{definition}{Definition}

\title{Compositional Algebra over Fixed-Width\\
Quaternionic-Symbolic State Packets}
\author{Cody Churchwell\\
\small Independent Researcher\\
\small \texttt{github.com/consigcody94/quaternion-monoid-algebra}}
\date{May 2026}

\begin{document}
\maketitle

\begin{abstract}
We define a compositional algebra over a fixed-width packet structure that
carries a unit quaternion, a small number of symbolic metadata fields, and a
positive scaling factor. A closed binary operation is defined on the packet
space with the property that the space, equipped with this operation, forms a
monoid: closure, two-sided identity, and associativity. The construction is
field-by-field; each sub-field operation is drawn from a small set of
associative monoid operations on its respective value space (Hamilton
multiplication on the quaternion factor; lookup-table substitution, lattice
maximum, modular addition, parity XOR, and multiplicative scaling on the
symbolic and scaling factors). The composite operation inherits associativity
from the per-field operations. We document an additional empirically-observed
property, topology preservation under iterated composition, and validate the
construction with a suite of fourteen automated tests including bit-exact
GPU/CPU correspondence and verification on a public motion-capture dataset.
Three application classes follow: composable agent-state evolution, algebraic
chain verification, and multi-source state composition. A permissively-licensed
reference implementation accompanies this paper.
\end{abstract}

\section{Introduction}

A recurring problem in robotics, swarm coordination, and accountability
infrastructure is the need for a bounded-width, time-evolving state
representation that is simultaneously composable, verifiable, and efficient
enough for edge deployment. Existing representations trade these properties
against one another. Raw floating-point quaternion streams are compact for
rotation but offer no composition operation beyond Hamilton multiplication on
the rotation component, and carry no symbolic metadata. Variable-width
cryptographic audit chains compose via Merkle structures but are not
fixed-width and are not amenable to single-cycle hardware composition.
High-dimensional vector embeddings compose via vector arithmetic but are lossy
and lack a unit-rotation guarantee.

This paper defines a fixed-width construction that fills the gap. A packet
carries a unit quaternion plus a small number of symbolic metadata fields plus
a scaling factor, with a defined composition operation whose result is itself a
packet of the same width. The resulting algebra has, simultaneously: closure
under composition, the algebraic structure of a monoid (so chains of operations
compose freely and associatively), and an empirically-observed
topology-preservation property under iterated composition.

\section{Construction}

\subsection{The packet structure}

A packet is the tuple
\[
p = (q,\; f_a,\; f_b,\; f_c,\; f_d,\; \pi,\; s),
\]
where $q \in S^3$ is a unit quaternion in $(w,x,y,z)$ layout; $f_a, f_b, f_c,
f_d$ are small symbolic fields of fixed bit width; $\pi \in \{0,1\}$ is a parity
bit; and $s \in \mathbb{R}_{>0}$ is a positive scaling factor. Field widths are
parameters of the construction. A representative default (widths $4,3,5,3$ bits,
plus parity, plus a smallest-three quaternion encoding in $32$ bits, plus an
IEEE-754 binary16 scaling factor) fits a $64$-bit packet.

\subsection{The composition operation}

For two packets $p_1$ and $p_2$, the composition $p_1 \otimes p_2$ is defined
field-by-field:

\begin{center}
\begin{tabular}{lll}
\toprule
Sub-field & Operation & Algebraic basis \\
\midrule
$q$    & Hamilton product, then normalize & $S^3$ group structure \\
$f_a$  & $T_a[i,j]$ lookup (default $i \oplus j$) & $\mathbb{Z}_2^{b}$ abelian group \\
$f_b$  & $\max(i,j)$ & totally-ordered lattice \\
$f_c$  & $(i+j) \bmod 2^{b_c}$ & $\mathbb{Z}/2^{b_c}\mathbb{Z}$ group \\
$f_d$  & $T_d[i,j]$ lookup (default $i \oplus j$) & $\mathbb{Z}_2^{b}$ abelian group \\
$\pi$  & $\pi_1 \oplus \pi_2$ & $\mathbb{Z}_2$ group \\
$s$    & $s_1 \cdot s_2$ & $\mathbb{R}_{>0}$ under multiplication \\
\bottomrule
\end{tabular}
\end{center}

The identity element is $I = (q_{\mathrm{id}}, 0, 0, 0, 0, 0, 1)$, where
$q_{\mathrm{id}} = (1,0,0,0)$.

\subsection{Monoid properties}

\begin{proposition}[Closure]
For all valid packets $p_1, p_2$, the product $p_1 \otimes p_2$ is a valid
packet.
\end{proposition}

\begin{proof}[Proof sketch]
Each per-field operation maps into the same field. The Hamilton product of two
unit quaternions, after normalization, is a unit quaternion. XOR and modular
addition remain in finite range. Lattice maximum of two values in
$[0, 2^b)$ remains in $[0,2^b)$. The product of two positive reals is a positive
real.
\end{proof}

\begin{proposition}[Two-sided identity]
$I \otimes p = p \otimes I = p$ for all valid $p$.
\end{proposition}

\begin{proof}[Proof sketch]
Field-wise: $q_{\mathrm{id}}$ is the quaternion identity; lookup tables
satisfy $T[0,j]=j$ and $T[i,0]=i$; $\max(0,x)=x$; $0+x \equiv x \pmod{N}$;
$0 \oplus x = x$; and $1 \cdot x = x$.
\end{proof}

\begin{proposition}[Associativity]
$(a \otimes b) \otimes c = a \otimes (b \otimes c)$ for all valid $a,b,c$.
\end{proposition}

\begin{proof}[Proof sketch]
Each per-field operation is associative on its value space: Hamilton product on
the quaternion group; XOR on $\mathbb{Z}_2^b$; maximum on a totally-ordered
lattice; modular addition on $\mathbb{Z}/N\mathbb{Z}$; parity XOR on
$\mathbb{Z}_2$; multiplication on $\mathbb{R}_{>0}$. A field-wise product of
associative operations is associative.
\end{proof}

Together these three propositions establish that the packet space, equipped
with $\otimes$, is a monoid.

\subsection{Topology preservation}

A property additional to the monoid structure: when the algebra is applied
iteratively to a stream of input packets via $\mathrm{state}[t{+}1] =
\mathrm{state}[t] \otimes \mathrm{stim}[t]$, the $H_1$ persistent-homology
signature of the resulting trajectory of quaternion components is bounded
relative to the $H_1$ signature of the input stream. Empirically, the ratio of
total bar persistence (output to input) falls in $[0.3, 5.0]$ for structured
inputs. We state this as an empirical observation rather than a theorem; a
formal characterization is left to future work.

\section{Validation}

The accompanying reference implementation passes fourteen automated tests:
eight algebraic-property tests (left and right identity over $100$ packets,
associativity over $512$ random triples, closure over $500$ random pairs,
stability over a $1000$-step self-product chain, power consistency, and
bit-exact GPU/CPU correspondence) and six stress tests ($10{,}000$-step
long-horizon stability, single-bit avalanche sensitivity, distinguishability
across $100$ chains, real-data behavior, field-saturation characterization, and
scale stability).

The GPU implementation agrees with the CPU reference bit-for-bit on the Hamilton
product: maximum absolute difference $0.00\mathrm{e}{+}00$ across $100{,}000$
random pairs at single precision. The real-data test verifies clean behavior on
the TUM RGB-D Pioneer 360 ground-truth sequence~\cite{Sturm2012}, downloaded at
runtime with SHA-256 verification, composing $2{,}000$ unit quaternions while
preserving the unit-norm constraint to machine precision.

The field-saturation test documents an honest limitation: the $\max(\cdot)$
sub-field saturates under iteration (it is associative with identity $0$ but is
not invertible), so it functions as a one-way ``high-amplitude event seen''
flag rather than a state register; the XOR and modular-addition fields are
information-preserving over their respective groups.

\section{Applications}

\paragraph{Composable agent-state evolution.}
$\mathrm{state}[t{+}1] = \mathrm{state}[t] \otimes \mathrm{stimulus}[t]$ is a
single bounded-width operation per time step. The agent's internal state is the
running composition of all stimuli received.

\paragraph{Algebraic chain verification.}
For an $N$-packet chain produced sequentially, a verifier computes the
cumulative product $I \otimes p_1 \otimes \cdots \otimes p_N$ and compares to a
separately-signed chain-head value. Verification runs in $O(N)$ operations at
constant per-operation cost. Empirically, a single-bit change anywhere in the
chain produces a divergent head, so the construction functions as a chain
digest under any cryptographic signature layered on top.

\paragraph{Multi-source state composition.}
$N$ independently-evolving state packets compose via the left-fold $p_1 \otimes
\cdots \otimes p_N$. By associativity the fold may be computed in any order or
as a parallel tree-reduce. By non-commutativity of the quaternion factor, the
order of sources is significant.

\section{Relationship to prior work}

The construction does not introduce new mathematics; each component has a long
literature: quaternion algebra~\cite{Hamilton1844}, persistent
homology~\cite{Edelsbrunner2002}, finite-group operations on $\mathbb{Z}_2^b$
and $\mathbb{Z}/N\mathbb{Z}$, and lookup-table substitution. The contribution
is the specific combination: a closed associative binary operation on a
fixed-width packet that simultaneously preserves a unit-quaternion integrity
check, symbolic-metadata semantics under per-field associative operations,
algebraic-chain composability, and an empirically-observed topology-preservation
property under iterated composition. We are not aware of this combination in a
single prior-art reference and welcome pointers to closer prior work.

\section{Availability}

A reference implementation (CPU and GPU), the full validation suite, worked
examples, and this paper are available under the MIT License at
\url{https://github.com/consigcody94/quaternion-monoid-algebra} and archived at
\href{https://doi.org/10.5281/zenodo.20301069}{doi:10.5281/zenodo.20301069}.

\paragraph{Provenance.}
This work was developed by the author with artificial-intelligence assistance
used as an engineering aid for implementation, validation, GPU porting, and
documentation. The conceptual direction and the specific algebraic structure
imposed on the packet space were the author's.

\begin{thebibliography}{9}

\bibitem{Hamilton1844}
W.~R. Hamilton.
\textit{On quaternions; or on a new system of imaginaries in algebra}.
Philosophical Magazine, 25(3):489--495, 1844.

\bibitem{Edelsbrunner2002}
H.~Edelsbrunner, D.~Letscher, and A.~Zomorodian.
\textit{Topological persistence and simplification}.
Discrete \& Computational Geometry, 28(4):511--533, 2002.

\bibitem{Sturm2012}
J.~Sturm, N.~Engelhard, F.~Endres, W.~Burgard, and D.~Cremers.
\textit{A benchmark for the evaluation of RGB-D SLAM systems}.
In Proc. International Conference on Intelligent Robot Systems (IROS), 2012.

\end{thebibliography}

\end{document}
